\setcounter{equation}{13}
\section{From balanced flows to typical-state mixing}\label{sec:mixing}

The flow estimate is available only when both sides of a cut have
appreciable Gibbs mass. The following finite-state argument is adapted
to exactly that information. It first finds a large set with internal
expansion, then uses the full chain's entropy dissipation to control
typical starting states. There is no requirement that a trajectory
remain in the large set.

\begin{lemma}[Pruning and averaged mixing]\label{lem:pruning}
Let $\Omega$ be a finite set, let $\pi$ be a strictly positive
probability on $\Omega$, and let $P_t$ be a continuous-time Markov
semigroup reversible with respect to $\pi$, with jump rates $c(x,y)$.
Assume that, for some $R>0$,
\[
 \sum_{y\ne x}c(x,y)\le R\qquad(x\in\Omega).
\]
Let $0<\varepsilon<1/3$ and $B>0$. Suppose
\[
 Q(S,S^{\rm c}):=\sum_{x\in S,\,y\notin S}\pi(x)c(x,y)>B
 \quad\text{whenever }\varepsilon\le\pi(S)\le1-\varepsilon.
\]
Then for every $T>0$,
\begin{equation}\label{eq:averaged-mixing}
 \sum_{x\in\Omega}\pi(x)\|P_T(x,\cdot)-\pi\|_{\rm TV}
 \le\frac{\sqrt{2R H(\pi)}}{B\sqrt T}+2\varepsilon,
 \qquad H(\pi)=-\sum_x\pi(x)\log\pi(x).
\end{equation}
\end{lemma}

\begin{proof}
Put $q(x,y)=\pi(x)c(x,y)=q(y,x)$, and use the unordered-edge
Dirichlet form
\[
 \mathfrak D(u)=\sum_{\{x,y\}}q(x,y)(u(x)-u(y))^2.
\]
Only distinct vertices occur, and every unordered pair is counted once.

\paragraph{Removing sets that do not expand.}
Start with $C=\Omega$. Whenever a nonempty $U\subset C$ satisfies
\[
 \pi(U)\le\pi(C)/2,\qquad Q(U,C\setminus U)<B\pi(U),
\]
remove $U$ from $C$. This process terminates because each removal
eliminates a vertex. If $A$ denotes all vertices removed so far,
then every edge between $A$ and its current complement was counted
when its first endpoint was removed. Thus, whenever $A\ne\varnothing$,
\[
 Q(A,\Omega\setminus A)<B\pi(A).
\]
Suppose the removed mass first reached $\varepsilon$. Its preceding
mass $a$ was less than $\varepsilon$, and the newly removed set has
mass at most $(1-a)/2$. The new total is therefore at most
$(1+a)/2<(1+\varepsilon)/2<1-\varepsilon$. This contradicts
the balanced-cut hypothesis, since $B\pi(A)\le B$.
Consequently the terminal set satisfies $\pi(C)>1-\varepsilon$
and
\begin{equation}\label{eq:internal-expansion}
 Q(U,C\setminus U)\ge B\pi(U)
 \quad\text{for every }U\subset C\text{ with }\pi(U)\le\pi(C)/2.
\end{equation}

\paragraph{A Poincar\'e inequality on the retained set.}
We use the coarea proof of the reversible Cheeger inequality
\cite[Theorem~3.5]{LawlerSokal1988}, keeping the original,
unnormalized weights on $C$. The resulting estimate is
\begin{equation}\label{eq:14}
 \sum_{y\in C}\pi(y)(u(y)-\bar u_C)^2
 \le\frac{2R}{B^2}\mathfrak D(u),
 \qquad \bar u_C=\frac1{\pi(C)}\sum_{y\in C}\pi(y)u(y).
\end{equation}
To prove it, choose a weighted median $m$ of $u$ on $C$. Each of
$w=(u-m)_+$ and $w=(m-u)_+$ has support of mass at most
$\pi(C)/2$. Integrating \eqref{eq:internal-expansion} over the
level sets of $w^2$ gives
\[
 B\sum_C\pi w^2
 \le\sum_{\{x,y\}\subset C}q(x,y)|w(x)^2-w(y)^2|.
\]
By Cauchy--Schwarz this is at most
\[
 \left(\sum_{\{x,y\}\subset C}q(x,y)(w(x)-w(y))^2\right)^{1/2}
 \left(2R\sum_C\pi w^2\right)^{1/2}.
\]
Indeed, $(w(x)+w(y))^2\le2(w(x)^2+w(y)^2)$, and the total outgoing
rate is at most $R$. If $\sum_C\pi w^2$ vanishes, its estimate
is immediate; otherwise divide and square. The sum of the energies
of the positive and negative parts is at most $\mathfrak D(u)$.
The weighted mean minimizes squared error, proving \eqref{eq:14}.

\paragraph{Entropy dissipation for the full chain.}
For a fixed starting state $x$, set $f_t^x(y)=P_t(x,y)/\pi(y)$.
Reversibility and
\[
 (a-b)(\log a-\log b)\ge4(\sqrt a-\sqrt b)^2
 \qquad(a,b>0)
\]
yield the entropy-dissipation estimate
\cite[Lemmas~2.5 and~2.7]{DiaconisSaloffCoste1996}
\begin{equation}\label{eq:15}
 \int_0^T\sum_x\pi(x)\mathfrak D(\sqrt{f_t^x})\,dt
 \le\frac14H(\pi).
\end{equation}
Indeed, the derivative of $\sum_y\pi(y)f_t^x(y)\log f_t^x(y)$ is
\[
 -\sum_{\{y,z\}}q(y,z)
       (f_t^x(y)-f_t^x(z))(\log f_t^x(y)-\log f_t^x(z)).
\]
The initial entropy is $\log(1/\pi(x))$, and the final entropy is
nonnegative. Averaging in $x\sim\pi$ and integrating proves
\eqref{eq:15}. Vanishing densities and reducible chains cause no
problem: start instead with $(1-\eta)f_0^x+\eta$ and pass to
$\eta\downarrow0$, or calculate on communicating classes.

\paragraph{Averaging fixed-state distances.}
Combining \eqref{eq:14} and \eqref{eq:15} supplies some
$t\in(0,T]$ for which
\[
 \sum_x\pi(x)r_x^2\le\frac{R H(\pi)}{2TB^2},
 \qquad r_x^2=\sum_C\pi(y)(\sqrt{f_t^x(y)}-a_x)^2,
 \qquad a_x=\frac1{\pi(C)}\sum_C\pi(y)\sqrt{f_t^x(y)}.
\]
Cauchy--Schwarz and
$a_x^2\pi(C)\le\sum_C\pi f_t^x\le1$ imply
\[
 \sum_C\pi|f_t^x-a_x^2|\le2r_x.
\]
Also, total normalization gives
$|\sum_C\pi(f_t^x-1)|\le\sum_{C^{\rm c}}\pi(f_t^x+1)$.
Comparing $f_t^x$ first with $a_x^2$ and then with $1$ on $C$,
and adding the complement, yields
\[
 \sum_\Omega\pi|f_t^x-1|
 \le4r_x+2\sum_{C^{\rm c}}\pi(f_t^x+1).
\]
The last sum before its factor two has average $2\pi(C^{\rm c})$
over $x\sim\pi$, because stationarity gives
$\sum_x\pi(x)f_t^x(y)=1$. Dividing by two for total variation
and applying Cauchy--Schwarz in $x$ proves the desired upper bound
at this time $t$. Contraction under $P_{T-t}$ gives the same bound
at the prescribed time $T$, establishing \eqref{eq:averaged-mixing}.
\end{proof}

\begin{proof}[Proof of Theorem~\ref{thm:main}]
On the disorder event of Proposition~\ref{prop:cuts}, apply
Lemma~\ref{lem:pruning} with
\[
 \varepsilon=n^{-d},\qquad B=\exp(-n^{1-b/20}),\qquad R=n.
\]
The last choice is valid because there are $n$ rate-one update
clocks. Also $H(\pi_G)\le n\log2$, since the state space has
$2^n$ elements. Hence
\[
 \sum_x\pi_G(x)\|P_T^G(x,\cdot)-\pi_G\|_{\rm TV}
 \le\frac{\sqrt{2\log2}\,n}{B\sqrt T}+2n^{-d}.
\]
Take $T=\exp(n^{1-b/40})$. The first term tends to zero, as
\[
 \log\frac{n}{B\sqrt T}
 =\log n+n^{1-b/20}-\tfrac12n^{1-b/40}\longrightarrow-\infty.
\]
The flow event itself has disorder probability tending to one.
Markov's inequality therefore shows that the Gibbs mass of fixed
states whose distance at time $T$ exceeds $1/4$ tends to zero in
disorder probability. The added diagonal changes neither this Gibbs
law nor the transition rates, so the same statement holds for the
model in \eqref{eq:intro-model}. This proves the theorem with
\[
 c_\beta=b/40=1/40000000.
\]
All threshold dimensions may depend on the fixed $\beta>1$.
\end{proof}
